\clearpage\begingroup\HistoricalHead{222}{L}{THEORETISCHE ASTRONOMIE. NACHLASS.}
\makeatletter\tagsleft@true\let\veqno\@@leqno\makeatother
\noindent{\small [S. 2]}
\begin{center}
[2.]\par\medskip
\textls[100]{Formeln zur ersten Annäherung aus dreien Beobachtungen.}
\end{center}
\begin{equation}
\begin{aligned}
\frac{R'}{\Delta'}\left(1-\frac{R^{\prime3}}{r^{\prime3}}\right)&={}\\
&\frac{\Theta\sin(\alpha''-\alpha')-\Theta'\sin(\alpha''-\alpha)+\Theta''\sin(\alpha'-\alpha)}
{\frac12(U'-U)(U''-U')\{\Theta''\sin(V'-\alpha)-\Theta\sin(V'-\alpha'')\}}
\text{[\textsuperscript{*})]},
\end{aligned}\tag*{I.}
\end{equation}
\begingroup\renewcommand{\thefootnote}{}\footnotetext{[\textsuperscript{*}) Vgl. Werke VI, S. 159.]}\endgroup
\begin{equation}
\begin{aligned}
\text{genau}\quad\left[\frac{1}{\Delta'R'R'}\left(1-\frac{R^{\prime3}}{r^{\prime3}}\right)\right]&={}\\
&\frac{d\alpha' dd\Theta'-d\Theta' dd\alpha'+\Theta'd\alpha^{\prime3}}
{dU^2\{\Theta'\cos(V'-\alpha')d\alpha'+\sin(V'-\alpha')d\Theta'\}}.
\end{aligned}\tag*{[a.]}
\end{equation}

Hat man nach dieser Formel $\dfrac{R'}{\Delta'}\left(1-\dfrac{R^{\prime3}}{r^{\prime3}}\right)$ bestimmt, so findet man darnach leicht durch Verbindung mit der Gleichung
\[
\frac{r'}{\Delta'}=\surd\left(1+\Theta'\Theta'+\frac{R'R'}{\Delta'\Delta'}-2\frac{R'}{\Delta'}\cos(V'-\alpha')\right)
\]
\noindent vermittelst weniger Versuche einen nahen Werth von $\Delta'$.
\begin{equation}
\begin{aligned}
\frac{\Delta''}{\Delta}
&=\frac{\Theta'\sin(V'-\alpha)-\Theta\sin(V'-\alpha')}
{\Theta''\sin(V'-\alpha')-\Theta'\sin(V'-\alpha'')}
\mathbin{.}\frac{t''-t'}{t'-t}\\
&=\frac{\dfrac{\sin(V'-\alpha)}{\sin(V'-\alpha')}-\dfrac{\Theta}{\Theta'}}
{\dfrac{\Theta''}{\Theta'}-\dfrac{\sin(V'-\alpha'')}{\sin(V'-\alpha')}}
\mathbin{.}\frac{t''-t'}{t'-t}\text{[\textsuperscript{**})]}.
\end{aligned}\tag*{II.}
\end{equation}
\begingroup\renewcommand{\thefootnote}{}\footnotetext{[\textsuperscript{**}) Vgl. Werke VI, S. 157.]}\endgroup
\noindent Völlig genau ist
\begin{equation}
\begin{aligned}
\frac{2d\Delta'}{\Delta'}&={}\\
&-\frac{\Theta'dd\alpha'\cos(V'-\alpha')+(\Theta'd\alpha^{\prime2}+dd\Theta')\sin(V'-\alpha')}
{\Theta'd\alpha'\cos(V'-\alpha')+d\Theta'\mathbin{.}\sin(V'-\alpha')}.
\end{aligned}\tag*{[b.]}
\end{equation}

In Ansehung der ersten Formel ist noch zu bemerken, dass $U'-U$,
\par\clearpage
\HistoricalHead{223}{R}{BESTIMMUNG DER BAHNEN DER HIMMELSKÖRPER.}\SourceLeaf{228}
\noindent $U''-U'$ in Theilen des Radius ausgedrückt werden müssen; auf die Weise ist $\log(\text{mot. med. }\odot\text{ diurn.})=8{,}2355820.792$\textup{[\textsuperscript{*})]}.
\begingroup\renewcommand{\thefootnote}{}\footnotetext{[\textsuperscript{*})] Wenn man auf die Erdmasse mit Rücksicht nimmt, muss man log mot. diurni Planetae in distantia media setzen: $8{,}2355814.21\mathbin{.}[A^{\frac32}:a^{\frac32}$, wo wie auch im Folgenden $A$ die halbe große Achse der Erdbahn bedeutet].}\endgroup

Nachdem man nun $\Delta'$ und $\dfrac{\Delta''}{\Delta}$ bestimmt hat, kann man, hinreichend genau zur ersten Näherung
\[
\begin{aligned}
\log\Delta&=\log\Delta'-\frac{t'-t}{t''-t}\log\frac{\Delta''}{\Delta},\\
\log\Delta''&=\log\Delta'+\frac{t''-t'}{t''-t}\log\frac{\Delta''}{\Delta}
\end{aligned}
\]
\noindent setzen.

Ist die Neigung der Bahn sehr gering, so sind obige Gleichungen unbrauchbar; die beobachtete Länge, ihre Veränderung und deren Zu- oder Abnahme gibt sodann bloss folgende Gleichung:
\begin{equation}
0=2\frac{d\Delta}{dU}\frac{d\alpha}{dU}+\Delta\frac{d^2\alpha}{dU^2}
+R\left(\frac{A^3}{R^3}-\frac{A^3}{r^3}\right)\sin(V-\alpha).
\tag*{[c.]}
\end{equation}
\noindent [Im übrigen ist:]
\begin{equation}
\begin{aligned}
0={}&\Theta\mathbin{.}\Delta\mathbin{.}\frac{d\alpha^2}{dU^2}
+2\frac{d\Delta}{dU}\frac{d\Theta}{dU}+\Delta\frac{d^2\Theta}{dU^2}\\
&-R\Theta\left(\frac{A^3}{R^3}-\frac{A^3}{r^3}\right)\cos(V-\alpha).
\end{aligned}\tag*{[d.]}
\end{equation}
\clearpage\endgroup
